% Standalone problem document, generated from the adjacent README.md.
% Compile with XeLaTeX. Mathematical content is maintained in Markdown.
\documentclass[11pt,a4paper]{article}
\usepackage[fontsize=10.5pt]{scrextend}
\usepackage[margin=22mm,headheight=15pt,headsep=9mm,footskip=11mm]{geometry}
\usepackage{fontspec}
\setmainfont{texgyrepagella}[Extension=.otf,UprightFont=*-regular,BoldFont=*-bold,ItalicFont=*-italic,BoldItalicFont=*-bolditalic]
\setsansfont{texgyreheros}[Extension=.otf,UprightFont=*-regular,BoldFont=*-bold,ItalicFont=*-italic,BoldItalicFont=*-bolditalic]
\setmonofont{texgyrecursor}[Extension=.otf,UprightFont=*-regular,BoldFont=*-bold,ItalicFont=*-italic,BoldItalicFont=*-bolditalic,Scale=MatchLowercase]
\usepackage{amsmath,amssymb}
\usepackage{unicode-math}
\setmathfont{texgyrepagella-math.otf}
\usepackage{xcolor,microtype,enumitem,needspace,fancyhdr}
\usepackage{bookmark}
\definecolor{ink}{HTML}{17344B}
\definecolor{accent}{HTML}{197D80}
\definecolor{muted}{HTML}{596773}
\hypersetup{colorlinks=true,linkcolor=accent,urlcolor=accent,pdftitle={IE-05 - Exact
extremizers for partial pivoting on orthogonal
matrices},pdfauthor={Open Problems in Numerical Linear Algebra}}
\urlstyle{same}
\setlength{\parindent}{0pt}
\setlength{\parskip}{5pt plus 1pt minus 1pt}
\linespread{1.04}
\setlength{\emergencystretch}{3em}
\widowpenalty=10000
\clubpenalty=10000
\displaywidowpenalty=10000
\allowdisplaybreaks[1]
\setlist{nosep,leftmargin=1.5em}
\providecommand{\tightlist}{\setlength{\itemsep}{0pt}\setlength{\parskip}{0pt}}
\providecommand{\pandocbounded}[1]{#1}
\pagestyle{fancy}
\fancyhf{}
\fancyhead[L]{\sffamily\footnotesize\color{muted}OPEN PROBLEMS IN NUMERICAL LINEAR ALGEBRA}
\fancyhead[R]{\sffamily\bfseries\footnotesize\color{accent}IE-05}
\fancyfoot[L]{\parbox[b]{0.9\textwidth}{\sffamily\fontsize{8}{10}\selectfont\color{muted}Editorial ratings; a bounded search cannot certify that a problem remains unsolved.\\Literature check: 2026-09-08}}
\fancyfoot[R]{\sffamily\footnotesize\color{muted}\thepage}
\renewcommand{\headrulewidth}{0.3pt}
\renewcommand{\headrule}{\hbox to\headwidth{\color{accent}\leaders\hrule height \headrulewidth\hfill}}
\makeatletter
\renewcommand{\section}{\Needspace{5\baselineskip}\@startsection{section}{1}{0pt}{12pt plus 2pt minus 2pt}{4pt}{\sffamily\bfseries\large\color{ink}}}
\renewcommand{\subsection}{\Needspace{4\baselineskip}\@startsection{subsection}{2}{0pt}{9pt plus 1pt minus 1pt}{3pt}{\sffamily\bfseries\normalsize\color{ink}}}
\makeatother
\setcounter{secnumdepth}{0}
\begin{document}
{\sffamily\small\color{muted}Linear systems and elimination\par}
\vspace{3pt}
{\raggedright\hyphenpenalty=10000\exhyphenpenalty=10000\sffamily\bfseries\fontsize{21}{24}\selectfont\color{ink}Exact
extremizers for partial pivoting on orthogonal matrices\par}
\vspace{7pt}
{\sffamily\small\textbf{Difficulty:} challenging\quad\textbf{Importance:} interesting
to specialist\par}
\vspace{4pt}
{\color{accent}\hrule height 0.8pt}
\vspace{6pt}
\textbf{Status:} open.

\subsection{Context and notation}\label{context-and-notation}

All elimination in this problem is in exact arithmetic. Write
\(\|A\|_{\max}=\max_{ij}|a_{ij}|\). A pivoting path creates successive
active Schur complements \(S_1=A,S_2,\ldots,S_n\), with row/column
permutations as appropriate. Its element-growth factor is

\[
\rho(A)=\frac{\max_{1\leq j\leq n}\|S_j\|_{\max}}{\|A\|_{\max}}.
\]

Partial pivoting chooses a largest-magnitude entry in the active first
column; complete pivoting chooses one anywhere in the active matrix. If
ties occur, a universal statement includes every admissible tie choice;
a supremum includes all admissible paths. These conventions remove
implementation-dependent ambiguity. Matrices are nonsingular unless
otherwise stated.

\subsection{Problem statement}\label{problem-statement}

Let \(L_n\) be the real unit lower triangular matrix whose entries
strictly below the diagonal are all \(-1\). Define \(Q_n\) by the unique
QR factorization \(L_n=Q_nR_n\) with positive diagonal in \(R_n\). For
\(Q_n\), use partial pivoting with the first available row chosen in a
tie. Is it true that, for every \(n\geq2\),

\[
\sup_{Q\in O(n)}\rho_{\mathrm{PP}}(Q)
=\rho_{\mathrm{PP}}(Q_n),
\qquad O(n)=\{Q\in\mathbb R^{n\times n}:Q^TQ=I\}?
\]

The supremum on the left includes all admissible partial-pivoting paths.
The candidate is fully specified by \(L_n\), rather than by an
approximate numerical optimizer. This asks for the sharp extremizer,
beyond the established exponential order of orthogonal growth.

\subsection{Reference}\label{reference}

Peca-Medlin, \href{https://arxiv.org/html/2308.16146v2}{\emph{Growth
factors of orthogonal matrices and local behavior of Gaussian
elimination with partial and complete pivoting}}, published in SIAM J.
Matrix Anal. Appl. (2024), §3.2 and Appendix B. The paper conjectures
this equality and establishes
\(\rho_{\mathrm{PP}}(Q_n)=2^{n-1}(1+o(1))/\sqrt3\).

\subsection{Status check}\label{status-check}

Searches for \textit{GEPP orthogonal conjecture 2026} and the exact
paper title found no proof of the extremal equality. The 2026 butterfly
paper still identifies orthogonal partial-pivoting growth as open. The
August Shah--Urschel results concern different growth questions and
pivot strategies; their exponential examples do not establish this exact
supremum.
\end{document}
